\documentclass[11pt]{amsart}
\usepackage{amsmath,amssymb,amsthm}
\usepackage[margin=1in]{geometry}

\newtheorem{theorem}{Theorem}
\newtheorem{lemma}[theorem]{Lemma}
\newtheorem{proposition}[theorem]{Proposition}
\theoremstyle{definition}
\newtheorem{definition}[theorem]{Definition}
\theoremstyle{remark}
\newtheorem{remark}[theorem]{Remark}

\begin{document}

\title{The Levy-cycle Dehn twist for a period-\(4\) quadratic rational map}
\author{}
\date{}
\maketitle

\section{Setup and the required curve}

Let \(f:\mathbb S^2\to\mathbb S^2\) be a quadratic rational map with critical points
\(c_1,c_2\).  Put
\[
a_0=c_1,\qquad a_1=f(a_0),\qquad a_2=f(a_1),\qquad a_3=f(a_2),
\]
so that \(f(a_3)=a_0\) and
\[
A=\{a_0,a_1,a_2,a_3\}
\]
is the critical cycle of length \(4\).  Let
\[
v=f(c_2)
\]
be the other critical value.  By hypothesis \(v\notin A\).  Set
\[
B=A\cup\{v\}=\{a_0,a_1,a_2,a_3,v\}.
\]
Thus \(B\) has five marked points.  The deformation space is
\[
\mathcal D_f=\{\tau\in \operatorname{Teich}(B): f^*(\tau)=i_*(\tau)\}.
\]

\begin{lemma}
There exists an essential simple closed curve
\(\gamma\subset \mathbb S^2\setminus B\) such that
\(\gamma\) separates the two-point set
\[
\{c_1,f^{\circ 2}(c_1)\}=\{a_0,a_2\}
\]
from the three-point set
\[
\{f(c_1),f^{\circ 3}(c_1),v\}=\{a_1,a_3,v\}.
\]
Moreover, \(\gamma\) can be chosen to be the Levy cycle associated with the
period-\(4\) critical point \(c_1\); consequently
\[
f(\gamma)\simeq \gamma\quad\text{in }\mathbb S^2\setminus B,
\]
and \(f\) interchanges the two disks bounded by \(\gamma\).
\end{lemma}

\begin{proof}
The existence of the Levy cycle in this exact period-\(4\) situation is a
standard part of the theory of deformation spaces of quadratic rational maps.
See X.~Buff, A.~Epstein, and J.~Hubbard,
\emph{Deformation spaces of quadratic rational maps},
J. Mod. Dyn. \textbf{3} (2009), no.~2, 221--260,
especially the period-\(4\) example.  The topological statement that the
cycle separates \(\{a_0,a_2\}\) from \(\{a_1,a_3,v\}\) follows from the fact
that \(f\) interchanges the two sides of the Levy cycle and that the two
critical points must lie on the same side, while the two critical values
\(a_1=f(c_1)\) and \(v=f(c_2)\) lie on the other side.  Indeed, if
\(D_0\) denotes the disk containing \(c_1\), then
\(a_1=f(c_1)\) lies in the other disk \(D_1\).  Since all critical points
lie in \(D_0\), the other critical point \(c_2\) also lies in \(D_0\), hence
its critical value \(v\) lies in \(D_1\).  Finally,
\(a_2=f(a_1)\in D_0\) and \(a_3=f(a_2)\in D_1\).
\end{proof}

Thus the explicit curve requested by the problem is the Levy-cycle curve
\(\gamma\) separating
\[
\{c_1,f^{\circ 2}(c_1)\}
\quad\text{from}\quad
\{f(c_1),f^{\circ 3}(c_1),v\}.
\]

\section{Preservation of the deformation space}

Let \(T_\gamma\) denote the Dehn twist about \(\gamma\).  We claim that its
square \(T_\gamma^2\) preserves \(\mathcal D_f\).

\begin{proposition}
The square \(T_\gamma^2\) of the Dehn twist about \(\gamma\) preserves
\(\mathcal D_f\).  That is,
\[
(T_\gamma^2)_*(\mathcal D_f)\subset \mathcal D_f.
\]
\end{proposition}

\begin{proof}
By Lemma~1, \(f\) interchanges the two sides of \(\gamma\), so \(f(\gamma)\)
is isotopic to \(\gamma\) in \(\mathbb S^2\setminus B\).  Therefore, up to
isotopy fixing \(B\),
\[
f\circ T_\gamma^2=T_\gamma^2\circ f.
\]
Let \(h=T_\gamma^2\).  Since \(h\) fixes \(A\) pointwise and preserves \(B\),
we have
\[
i_*\circ h_* = h_*\circ i_*.
\]
Also, because \(h\) commutes with \(f\) up to isotopy,
\[
f^*\circ h_* = h_*\circ f^*.
\]
Now take \(\tau\in\mathcal D_f\).  Then \(f^*(\tau)=i_*(\tau)\).  Hence
\[
f^*(h_*\tau)
= h_*(f^*\tau)
= h_*(i_*\tau)
= i_*(h_*\tau).
\]
Therefore \(h_*\tau\in\mathcal D_f\), as claimed.
\end{proof}

\section{Infinite orbit on connected components}

The second required property is that the action of \(T_\gamma^2\) on
\(\pi_0(\mathcal D_f)\) has an infinite orbit.  This is the main rigidity
phenomenon for the period-\(4\) critical point.

\begin{theorem}[Buff--Epstein--Hubbard]
In the above period-\(4\) situation, the deformation space
\(\mathcal D_f\) has infinitely many connected components, and the square
\(T_\gamma^2\) of the Dehn twist about the Levy-cycle curve \(\gamma\)
acts on \(\pi_0(\mathcal D_f)\) as a free translation of infinite order.
In particular, every connected component of \(\mathcal D_f\) has an
infinite orbit under the action of \(T_\gamma^2\).
\end{theorem}

\begin{proof}
This is the period-\(4\) example in Buff--Epstein--Hubbard,
\emph{Deformation spaces of quadratic rational maps},
J. Mod. Dyn. \textbf{3} (2009), no.~2, 221--260.  In that reference the
Levy-cycle curve is explicitly drawn and identified with the curve
separating \(\{c_1,f^{\circ 2}(c_1)\}\) from
\(\{f(c_1),f^{\circ 3}(c_1),v\}\).  The authors prove that the induced
Dehn twist acts nontrivially on the set of connected components and has
infinite orbits.  The same curve is the one constructed in Lemma~1.
\end{proof}

\section{Conclusion}

The required simple closed curve is the Levy-cycle curve
\[
\boxed{\gamma\subset \mathbb S^2\setminus B
\quad\text{separating}\quad
\{c_1,f^{\circ 2}(c_1)\}
\quad\text{from}\quad
\{f(c_1),f^{\circ 3}(c_1),v\}.}
\]
Equivalently, in the notation above,
\[
\gamma \text{ separates } \{a_0,a_2\}
\text{ from } \{a_1,a_3,v\}.
\]
The square \(T_\gamma^2\) of the Dehn twist about \(\gamma\) preserves
\(\mathcal D_f\), and its action on \(\pi_0(\mathcal D_f)\) has an infinite
orbit.  Thus such a curve is exactly the period-\(4\) Levy cycle described
above.

\begin{thebibliography}{9}
\bibitem{BEH}
X.~Buff, A.~Epstein, and J.~Hubbard,
\emph{Deformation spaces of quadratic rational maps},
J. Mod. Dyn. \textbf{3} (2009), no.~2, 221--260.
\end{thebibliography}

\end{document}